\chapter{Hydrodynamics from Ward Identities}
\label{ch:thermal-hydrodynamics-from-ward-identities}

\ConventionsLorentzian

Hydrodynamics is the long-wavelength effective theory of conserved densities.
It is not obtained by assuming microscopic quasiparticles.  The input is a
thermal or locally thermal state, Ward identities for conserved currents, and
a derivative expansion of expectation values in terms of slowly varying
intensive variables.

\section{Conserved Densities}

Let \(T^{\mu\nu}\) be the stress tensor and \(J^\mu_A\) conserved currents.
In flat space without external sources,
\[
  \partial_\mu T^{\mu\nu}=0,\qquad
  \partial_\mu J^\mu_A=0 .
\]
In a state near local equilibrium, introduce a velocity field \(u^\mu(x)\)
with
\[
  u^\mu u_\mu=-1,
\]
temperature \(T(x)\), and chemical potentials \(\mu_A(x)\).  The ideal
constitutive relations are
\[
  T^{\mu\nu}_{(0)}
  =
  \varepsilon u^\mu u^\nu
  +p\,\Delta^{\mu\nu},
  \qquad
  J^\mu_{A(0)}=n_A u^\mu,
\]
where
\[
  \Delta^{\mu\nu}=\eta^{\mu\nu}+u^\mu u^\nu
\]
projects orthogonally to the velocity.

\section{Thermodynamic Closure}

The functions \(\varepsilon(T,\mu)\), \(p(T,\mu)\), and \(n_A(T,\mu)\) are
related by the local first law
\[
  d\varepsilon=T\,ds+\mu_A\,dn_A,
  \qquad
  dp=s\,dT+n_A\,d\mu_A,
\]
and
\[
  \varepsilon+p=Ts+\mu_A n_A .
\]
These equations close the ideal hydrodynamic system after an equation of
state is specified by the underlying QFT.

\section{Entropy Current at Ideal Order}

Define
\[
  S^\mu_{(0)}=s u^\mu .
\]
Using the Ward identities and the thermodynamic relations, one finds
\[
  \partial_\mu S^\mu_{(0)}=0
\]
for ideal nondissipative flow.  Indeed,
\[
  u_\nu\partial_\mu T^{\mu\nu}=0
\]
gives
\[
  u^\mu\partial_\mu\varepsilon
  +(\varepsilon+p)\partial_\mu u^\mu=0,
\]
while current conservation gives
\[
  u^\mu\partial_\mu n_A+n_A\partial_\mu u^\mu=0 .
\]
Substituting the first law cancels all terms in
\(\partial_\mu(su^\mu)\).

\section{First-Order Dissipation}

At first derivative order in a parity-even relativistic fluid,
\[
  T^{\mu\nu}
  =
  T^{\mu\nu}_{(0)}
  -\eta\,\sigma^{\mu\nu}
  -\zeta\,\Delta^{\mu\nu}\partial_\lambda u^\lambda+\cdots ,
\]
where
\[
  \sigma^{\mu\nu}
  =
  2\Delta^{\mu\alpha}\Delta^{\nu\beta}
  \left(
  \partial_{(\alpha}u_{\beta)}
  -\frac{1}{D-1}\eta_{\alpha\beta}\partial_\lambda u^\lambda
  \right).
\]
Charge diffusion is represented at first derivative order by
\[
  J^\mu_A=n_Au^\mu-\Sigma_{AB}\Delta^{\mu\nu}
  \partial_\nu\!\left(\frac{\mu_B}{T}\right)+\cdots .
\]
The local second law requires
\[
  \eta\geq0,\qquad \zeta\geq0,\qquad
  \Sigma_{AB}\ \text{positive semidefinite}.
\]

\section{Linearized Modes}

Linearize a neutral fluid around rest:
\[
  u^\mu=(1,\mathbf v),\qquad
  T=T_0+\delta T .
\]
The transverse velocity obeys
\[
  \partial_t v_\perp-\frac{\eta}{\varepsilon+p}\nabla^2v_\perp=0,
\]
so the shear pole is
\[
  \omega=-\ii\,\frac{\eta}{\varepsilon+p}k^2+\cdots .
\]
Longitudinal fluctuations give sound poles
\[
  \omega=\pm c_s k-\frac{\ii}{2}\Gamma_s k^2+\cdots ,
  \qquad
  c_s^2=\left(\frac{\partial p}{\partial\varepsilon}\right)_{s/n},
\]
with attenuation \(\Gamma_s\) determined by viscosities and charge diffusion
data.  These poles are the analytic continuation of the retarded correlators
whose zero-frequency slopes define transport in
Chapter~\ref{ch:thermal-spectral-kubo-transport}.

\section{Hydrodynamic Limit as a QFT Claim}

A QFT derivation of hydrodynamics requires more than conservation laws.  One
must prove that nonconserved excitations decay on microscopic time scales,
that the derivative expansion is asymptotic in a controlled regime, and that
the Kubo coefficients are the same coefficients appearing in the constitutive
relations.

\begin{openproblem}[Hydrodynamic emergence theorem]
\label{op:hydrodynamic-emergence-theorem}
Find physically useful QFT hypotheses under which local thermalization,
clustering, analyticity of retarded functions, positivity of entropy
production, and the hydrodynamic pole expansion can be proved from the
microscopic theory.
\end{openproblem}
